\section{What this suggests}

The central question is not only how accurately a tokenizer reconstructs
actions, but also how easily a single autoregressive transformer can predict its
tokens. A representation with low quantization error can still be difficult to
generate if it requires long sequences or concentrates predictive difficulty in
a few token positions.

This suggests that action tokenization should be designed together with the
autoregressive policy. Compression can shorten the generation horizon, while a
more even allocation of information may make token prediction easier. Both goals
must preserve the action detail needed for control.

The evaluations on RoboMimic, RoboCasa, and real-world tasks will test whether
these design choices translate into better policies. Improvements across these
settings would support the view that tokenization is a central part of making
the single-transformer autoregressive VLA setting work.
